% Appendix: Supplementary Material
% Included via % Appendix: Supplementary Material
% Included via % Appendix: Supplementary Material
% Included via % Appendix: Supplementary Material
% Included via \include{appendix} from report.tex (after \appendix command)

\chapter{Supplementary Material}
\label{ch:appendix}

%=================================================================
\section{Full Per-Fold Results: All Seven Encodings}
\label{sec:app-encoding-results}
%=================================================================

Table~\ref{tab:full-encoding} reports the complete per-fold accuracy for all seven spike encoding methods and the ANN baseline. All experiments use the SpikingCNN architecture ($\sim$622K parameters), ESC-50 five-fold cross-validation (400 test samples per fold), Adam optimiser (lr $= 10^{-3}$, weight decay $= 10^{-4}$), 50 epochs maximum with early stopping at patience $= 10$, and $T = 25$ timesteps. The ANN uses the identical architecture with ReLU activations and no temporal dimension. Population encoding uses 500 output neurons (10 per class) with MSE count loss; all other SNN variants use cross-entropy on summed membrane potentials.

\begin{table}[H]
\centering
\caption{Complete per-fold accuracy (\%) for all encodings on ESC-50 five-fold cross-validation.}
\label{tab:full-encoding}
\small
\begin{tabular}{lccccccr}
\toprule
\textbf{Encoding} & \textbf{Fold 1} & \textbf{Fold 2} & \textbf{Fold 3} & \textbf{Fold 4} & \textbf{Fold 5} & \textbf{Mean} & \textbf{Std} \\
\midrule
ANN (baseline) & 63.25 & 59.50 & 65.25 & 68.75 & 62.50 & 63.85 & 3.07 \\
\midrule
SNN Direct     & 40.50 & 48.50 & 48.25 & 54.00 & 44.50 & 47.15 & 4.50 \\
SNN Rate       & 24.50 & 27.25 & 23.00 & 21.50 & 23.75 & 24.00 & 1.90 \\
SNN Phase      & 22.50 & 22.25 & 25.00 & 24.25 & 26.75 & 24.15 & 1.66 \\
SNN Population & 22.75 & 18.50 & 15.75 & 22.00 & 16.75 & 19.15 & 2.79 \\
SNN Latency    & 14.00 & 15.75 & 17.75 & 15.50 & 18.50 & 16.30 & 1.62 \\
SNN Delta      &  8.25 &  7.75 &  7.25 &  7.50 &  5.50 &  7.25 & 0.94 \\
SNN Burst      &  5.00 &  5.25 &  9.25 &  6.00 &  7.00 &  6.50 & 1.54 \\
\bottomrule
\end{tabular}
\end{table}

%=================================================================
\section{PANNs Transfer Learning: Per-Fold Results}
\label{sec:app-panns}
%=================================================================

Table~\ref{tab:panns-folds} reports the per-fold accuracy for the PANNs transfer learning experiment. Frozen CNN14 embeddings (2,048 dimensions, pretrained on AudioSet) are classified by three heads: a three-layer SNN, a matched three-layer ANN (ReLU), and a linear logistic regression baseline. All heads are trained for 50 epochs with the same Adam optimiser and learning rate schedule.

\begin{table}[H]
\centering
\caption{PANNs transfer learning: per-fold accuracy (\%) on ESC-50.}
\label{tab:panns-folds}
\small
\begin{tabular}{lccccccc}
\toprule
\textbf{Model} & \textbf{Fold 1} & \textbf{Fold 2} & \textbf{Fold 3} & \textbf{Fold 4} & \textbf{Fold 5} & \textbf{Mean} & \textbf{Std} \\
\midrule
PANNs + SNN    & 92.00 & 94.50 & 91.00 & 93.50 & 91.50 & 92.50 & 1.30 \\
PANNs + ANN    & 93.00 & 95.00 & 92.00 & 95.50 & 91.75 & 93.45 & 1.54 \\
PANNs + Linear & 94.25 & 95.75 & 92.50 & 95.25 & 91.25 & 93.80 & 1.69 \\
\bottomrule
\end{tabular}
\end{table}

%=================================================================
\section{SpiNNaker Calibration Parameters}
\label{sec:app-spinnaker-params}
%=================================================================

Table~\ref{tab:spinnaker-neuron} lists the neuron model parameters used for all SpiNNaker deployments. These were determined through a systematic parameter sweep on a 20-sample pilot set (Run~5).

\begin{table}[H]
\centering
\caption{SpiNNaker neuron model parameters (IF\_curr\_exp).}
\label{tab:spinnaker-neuron}
\small
\begin{tabular}{llp{8cm}}
\toprule
\textbf{Parameter} & \textbf{Value} & \textbf{Rationale} \\
\midrule
Neuron model       & IF\_curr\_exp       & sPyNNaker default; approximates snnTorch LIF \\
$v_\text{thresh}$  & 1.0\,mV            & Matches snnTorch threshold $\vartheta = 1.0$ \\
$v_\text{rest}$    & 0.0\,mV            & Matches snnTorch reset-to-zero semantics \\
$v_\text{reset}$   & 0.0\,mV            & Hard reset after spike emission \\
$\tau_m$           & 20.0\,ms           & Calibrated via scale sweep; corresponds to $\beta \approx 0.95$ \\
$\tau_\text{syn,E}$& 5.0\,ms            & Faster decay than $\tau_m$ for temporal precision \\
$\tau_\text{syn,I}$& 5.0\,ms            & Symmetric (no inhibitory synapses used) \\
$C_m$              & 1.0\,nF            & Default capacitance \\
$I_\text{offset}$  & 0.0\,nA            & No bias current injected \\
\bottomrule
\end{tabular}
\end{table}

Table~\ref{tab:spinnaker-board} summarises the SpiNNaker board configuration and simulation parameters.

\begin{table}[H]
\centering
\caption{SpiNNaker board and simulation configuration.}
\label{tab:spinnaker-board}
\small
\begin{tabular}{ll}
\toprule
\textbf{Property} & \textbf{Value} \\
\midrule
Board              & SpiNN-5 (48 chips, 864 ARM cores) \\
Access             & \texttt{spinnaker.cs.man.ac.uk} via sPyNNaker \\
sPyNNaker version  & 1.0.0 \\
Simulation timestep& 1.0\,ms \\
Duration per sample& 25\,ms ($T = 25$ timesteps) \\
FC2 layer size     & 256 input neurons $\to$ 50 output neurons \\
Connectivity       & FromListConnector (pre-computed weight matrix) \\
Recording          & Membrane potential at each timestep for all 50 output neurons \\
\bottomrule
\end{tabular}
\end{table}

%=================================================================
\section{SpiNNaker Pruning Results}
\label{sec:app-spinnaker-pruning}
%=================================================================

Table~\ref{tab:spinnaker-pruning} reports the mean SpiNNaker and snnTorch accuracy across five folds at each pruning level. Pruning is applied as unstructured magnitude pruning to the FC2 weight matrix; weights below the $k$-th percentile are set to zero.

\begin{table}[H]
\centering
\caption{SpiNNaker and snnTorch accuracy (\%) by pruning level (five-fold means).}
\label{tab:spinnaker-pruning}
\small
\begin{tabular}{lccc}
\toprule
\textbf{Pruning (\%)} & \textbf{SpiNNaker} & \textbf{snnTorch} & \textbf{Gap (pp)} \\
\midrule
0  (baseline)  & 33.1 & 46.0 & 12.9 \\
50             & 33.8 & 44.2 & 10.4 \\
55             & 34.1 & 43.5 &  9.4 \\
60             & 35.2 & 42.8 &  7.6 \\
65             & 34.6 & 41.3 &  6.7 \\
70             & 33.9 & 39.7 &  5.8 \\
75             & 32.4 & 37.1 &  4.7 \\
80             & 30.8 & 33.5 &  2.7 \\
85             & 26.1 & 28.9 &  2.8 \\
90             & 19.5 & 21.0 &  1.5 \\
95             &  8.3 &  9.2 &  0.9 \\
\bottomrule
\end{tabular}
\end{table}

The software--hardware gap narrows monotonically with pruning. At 60\% pruning, the SpiNNaker accuracy (35.2\%) is higher than the unpruned software baseline might suggest given the gap trend. At 80\% pruning, the gap has narrowed to 2.7 pp. This supports the interpretation that small-magnitude weights---those most susceptible to fixed-point quantisation error---are disproportionately responsible for the hardware accuracy deficit.

%=================================================================
\section{Reproducibility}
\label{sec:app-reproducibility}
%=================================================================

\subsection{Random Seeds}

Table~\ref{tab:seeds} documents the random seeds used across all experiments to ensure reproducibility.

\begin{table}[H]
\centering
\caption{Random seed configuration for all experiments.}
\label{tab:seeds}
\small
\begin{tabular}{lll}
\toprule
\textbf{Experiment} & \textbf{Seed} & \textbf{Notes} \\
\midrule
All 5-fold training        & 42 (torch, numpy)  & Set at start of each fold \\
CSF3 cluster training      & fold number (1--5) & Per-fold in SLURM submission \\
Surrogate ablation         & 42                 & Fixed single seed, fold 1 \\
t-SNE visualisation        & 42                 & \texttt{sklearn.TSNE random\_state} \\
PANNs + SNN head           & 42                 & All 5 folds \\
Adversarial attacks        & N/A                & Deterministic given model weights \\
Continual learning         & 42                 & Pretrained from fold 4 best \\
\bottomrule
\end{tabular}
\end{table}

\subsection{Software Versions}

Table~\ref{tab:software} lists all software dependencies and their versions.

\begin{table}[H]
\centering
\caption{Software dependencies.}
\label{tab:software}
\small
\begin{tabular}{lll}
\toprule
\textbf{Package} & \textbf{Version} & \textbf{Purpose} \\
\midrule
Python           & 3.14             & Runtime \\
PyTorch          & 2.10             & Deep learning framework \\
snnTorch         & 0.9.4            & SNN neuron models and training \\
librosa          & 0.10.x           & Audio loading and mel spectrogram computation \\
NeuroBench       & 2.2.0            & Energy and synaptic operation metrics \\
panns-inference  & latest           & CNN14 embedding extraction \\
torchattacks     & latest           & FGSM and PGD adversarial attacks \\
sPyNNaker        & 1.0.0            & SpiNNaker deployment (Python 3.11 venv) \\
scikit-learn     & latest           & t-SNE, classification metrics \\
scipy            & latest           & Wilcoxon signed-rank test \\
\bottomrule
\end{tabular}
\end{table}

\subsection{Hardware}

\begin{itemize}[nosep]
    \item \textbf{GPU training:} NVIDIA A100-SXM4-80GB on the University of Manchester CSF3 cluster. CUDA 12.6.2. Maximum two GPUs per user.
    \item \textbf{Local development:} Apple M-series with MPS backend (PyTorch). Used for code development and debugging only; all reported results are from CSF3 A100 runs.
    \item \textbf{SpiNNaker:} SpiNN-5 board (48 SpiNNaker chips, 864 ARM968 cores) accessed remotely via \texttt{spinnaker.cs.man.ac.uk}. sPyNNaker 1.0.0.
\end{itemize}

\subsection{Data Preprocessing Pipeline}

The exact preprocessing pipeline, implemented in \texttt{src/dataset.py}, is as follows:

\begin{lstlisting}[language=Python, caption={Audio preprocessing pipeline.}]
y, sr = librosa.load(filepath, sr=22050, duration=5)
expected_len = 22050 * 5  # 110,250 samples
if len(y) < expected_len:
    y = np.pad(y, (0, expected_len - len(y)))
mel = librosa.feature.melspectrogram(
    y=y, sr=sr, n_mels=64, n_fft=1024, hop_length=512
)
mel_db = librosa.power_to_db(mel, ref=np.max)
mel_norm = (mel_db - mel_db.min()) / (mel_db.max() - mel_db.min() + 1e-8)
# Final shape: (1, 64, 216)
tensor = torch.tensor(mel_norm, dtype=torch.float32).unsqueeze(0)
\end{lstlisting}

\subsection{Dataset}

ESC-50 is publicly available at \url{https://github.com/karolpiczak/ESC-50}. It comprises 2{,}000 five-second recordings across 50 classes with predefined five-fold cross-validation splits. No modifications were made to the dataset. Audio files are automatically downloaded and cached at \texttt{data/ESC-50-master/} on first execution.

\subsection{Code Availability}

All source code, training scripts, experiment scripts, analysis notebooks, and result JSON files are available in the project repository. The repository structure is:

\begin{lstlisting}[basicstyle=\small\ttfamily, frame=none]
snn-esc50/
  src/models/        -- SpikingCNN, ConvANN, encoding modules
  src/dataset.py     -- ESC-50 data loading with librosa
  src/train.py       -- 5-fold CV training entry point
  experiments/       -- All advanced experiment scripts
  spinnaker/         -- SpiNNaker deployment scripts
  results/           -- JSON result files (model checkpoints excluded)
\end{lstlisting}
 from report.tex (after \appendix command)

\chapter{Supplementary Material}
\label{ch:appendix}

%=================================================================
\section{Full Per-Fold Results: All Seven Encodings}
\label{sec:app-encoding-results}
%=================================================================

Table~\ref{tab:full-encoding} reports the complete per-fold accuracy for all seven spike encoding methods and the ANN baseline. All experiments use the SpikingCNN architecture ($\sim$622K parameters), ESC-50 five-fold cross-validation (400 test samples per fold), Adam optimiser (lr $= 10^{-3}$, weight decay $= 10^{-4}$), 50 epochs maximum with early stopping at patience $= 10$, and $T = 25$ timesteps. The ANN uses the identical architecture with ReLU activations and no temporal dimension. Population encoding uses 500 output neurons (10 per class) with MSE count loss; all other SNN variants use cross-entropy on summed membrane potentials.

\begin{table}[H]
\centering
\caption{Complete per-fold accuracy (\%) for all encodings on ESC-50 five-fold cross-validation.}
\label{tab:full-encoding}
\small
\begin{tabular}{lccccccr}
\toprule
\textbf{Encoding} & \textbf{Fold 1} & \textbf{Fold 2} & \textbf{Fold 3} & \textbf{Fold 4} & \textbf{Fold 5} & \textbf{Mean} & \textbf{Std} \\
\midrule
ANN (baseline) & 63.25 & 59.50 & 65.25 & 68.75 & 62.50 & 63.85 & 3.07 \\
\midrule
SNN Direct     & 40.50 & 48.50 & 48.25 & 54.00 & 44.50 & 47.15 & 4.50 \\
SNN Rate       & 24.50 & 27.25 & 23.00 & 21.50 & 23.75 & 24.00 & 1.90 \\
SNN Phase      & 22.50 & 22.25 & 25.00 & 24.25 & 26.75 & 24.15 & 1.66 \\
SNN Population & 22.75 & 18.50 & 15.75 & 22.00 & 16.75 & 19.15 & 2.79 \\
SNN Latency    & 14.00 & 15.75 & 17.75 & 15.50 & 18.50 & 16.30 & 1.62 \\
SNN Delta      &  8.25 &  7.75 &  7.25 &  7.50 &  5.50 &  7.25 & 0.94 \\
SNN Burst      &  5.00 &  5.25 &  9.25 &  6.00 &  7.00 &  6.50 & 1.54 \\
\bottomrule
\end{tabular}
\end{table}

%=================================================================
\section{PANNs Transfer Learning: Per-Fold Results}
\label{sec:app-panns}
%=================================================================

Table~\ref{tab:panns-folds} reports the per-fold accuracy for the PANNs transfer learning experiment. Frozen CNN14 embeddings (2,048 dimensions, pretrained on AudioSet) are classified by three heads: a three-layer SNN, a matched three-layer ANN (ReLU), and a linear logistic regression baseline. All heads are trained for 50 epochs with the same Adam optimiser and learning rate schedule.

\begin{table}[H]
\centering
\caption{PANNs transfer learning: per-fold accuracy (\%) on ESC-50.}
\label{tab:panns-folds}
\small
\begin{tabular}{lccccccc}
\toprule
\textbf{Model} & \textbf{Fold 1} & \textbf{Fold 2} & \textbf{Fold 3} & \textbf{Fold 4} & \textbf{Fold 5} & \textbf{Mean} & \textbf{Std} \\
\midrule
PANNs + SNN    & 92.00 & 94.50 & 91.00 & 93.50 & 91.50 & 92.50 & 1.30 \\
PANNs + ANN    & 93.00 & 95.00 & 92.00 & 95.50 & 91.75 & 93.45 & 1.54 \\
PANNs + Linear & 94.25 & 95.75 & 92.50 & 95.25 & 91.25 & 93.80 & 1.69 \\
\bottomrule
\end{tabular}
\end{table}

%=================================================================
\section{SpiNNaker Calibration Parameters}
\label{sec:app-spinnaker-params}
%=================================================================

Table~\ref{tab:spinnaker-neuron} lists the neuron model parameters used for all SpiNNaker deployments. These were determined through a systematic parameter sweep on a 20-sample pilot set (Run~5).

\begin{table}[H]
\centering
\caption{SpiNNaker neuron model parameters (IF\_curr\_exp).}
\label{tab:spinnaker-neuron}
\small
\begin{tabular}{llp{8cm}}
\toprule
\textbf{Parameter} & \textbf{Value} & \textbf{Rationale} \\
\midrule
Neuron model       & IF\_curr\_exp       & sPyNNaker default; approximates snnTorch LIF \\
$v_\text{thresh}$  & 1.0\,mV            & Matches snnTorch threshold $\vartheta = 1.0$ \\
$v_\text{rest}$    & 0.0\,mV            & Matches snnTorch reset-to-zero semantics \\
$v_\text{reset}$   & 0.0\,mV            & Hard reset after spike emission \\
$\tau_m$           & 20.0\,ms           & Calibrated via scale sweep; corresponds to $\beta \approx 0.95$ \\
$\tau_\text{syn,E}$& 5.0\,ms            & Faster decay than $\tau_m$ for temporal precision \\
$\tau_\text{syn,I}$& 5.0\,ms            & Symmetric (no inhibitory synapses used) \\
$C_m$              & 1.0\,nF            & Default capacitance \\
$I_\text{offset}$  & 0.0\,nA            & No bias current injected \\
\bottomrule
\end{tabular}
\end{table}

Table~\ref{tab:spinnaker-board} summarises the SpiNNaker board configuration and simulation parameters.

\begin{table}[H]
\centering
\caption{SpiNNaker board and simulation configuration.}
\label{tab:spinnaker-board}
\small
\begin{tabular}{ll}
\toprule
\textbf{Property} & \textbf{Value} \\
\midrule
Board              & SpiNN-5 (48 chips, 864 ARM cores) \\
Access             & \texttt{spinnaker.cs.man.ac.uk} via sPyNNaker \\
sPyNNaker version  & 1.0.0 \\
Simulation timestep& 1.0\,ms \\
Duration per sample& 25\,ms ($T = 25$ timesteps) \\
FC2 layer size     & 256 input neurons $\to$ 50 output neurons \\
Connectivity       & FromListConnector (pre-computed weight matrix) \\
Recording          & Membrane potential at each timestep for all 50 output neurons \\
\bottomrule
\end{tabular}
\end{table}

%=================================================================
\section{SpiNNaker Pruning Results}
\label{sec:app-spinnaker-pruning}
%=================================================================

Table~\ref{tab:spinnaker-pruning} reports the mean SpiNNaker and snnTorch accuracy across five folds at each pruning level. Pruning is applied as unstructured magnitude pruning to the FC2 weight matrix; weights below the $k$-th percentile are set to zero.

\begin{table}[H]
\centering
\caption{SpiNNaker and snnTorch accuracy (\%) by pruning level (five-fold means).}
\label{tab:spinnaker-pruning}
\small
\begin{tabular}{lccc}
\toprule
\textbf{Pruning (\%)} & \textbf{SpiNNaker} & \textbf{snnTorch} & \textbf{Gap (pp)} \\
\midrule
0  (baseline)  & 33.1 & 46.0 & 12.9 \\
50             & 33.8 & 44.2 & 10.4 \\
55             & 34.1 & 43.5 &  9.4 \\
60             & 35.2 & 42.8 &  7.6 \\
65             & 34.6 & 41.3 &  6.7 \\
70             & 33.9 & 39.7 &  5.8 \\
75             & 32.4 & 37.1 &  4.7 \\
80             & 30.8 & 33.5 &  2.7 \\
85             & 26.1 & 28.9 &  2.8 \\
90             & 19.5 & 21.0 &  1.5 \\
95             &  8.3 &  9.2 &  0.9 \\
\bottomrule
\end{tabular}
\end{table}

The software--hardware gap narrows monotonically with pruning. At 60\% pruning, the SpiNNaker accuracy (35.2\%) is higher than the unpruned software baseline might suggest given the gap trend. At 80\% pruning, the gap has narrowed to 2.7 pp. This supports the interpretation that small-magnitude weights---those most susceptible to fixed-point quantisation error---are disproportionately responsible for the hardware accuracy deficit.

%=================================================================
\section{Reproducibility}
\label{sec:app-reproducibility}
%=================================================================

\subsection{Random Seeds}

Table~\ref{tab:seeds} documents the random seeds used across all experiments to ensure reproducibility.

\begin{table}[H]
\centering
\caption{Random seed configuration for all experiments.}
\label{tab:seeds}
\small
\begin{tabular}{lll}
\toprule
\textbf{Experiment} & \textbf{Seed} & \textbf{Notes} \\
\midrule
All 5-fold training        & 42 (torch, numpy)  & Set at start of each fold \\
CSF3 cluster training      & fold number (1--5) & Per-fold in SLURM submission \\
Surrogate ablation         & 42                 & Fixed single seed, fold 1 \\
t-SNE visualisation        & 42                 & \texttt{sklearn.TSNE random\_state} \\
PANNs + SNN head           & 42                 & All 5 folds \\
Adversarial attacks        & N/A                & Deterministic given model weights \\
Continual learning         & 42                 & Pretrained from fold 4 best \\
\bottomrule
\end{tabular}
\end{table}

\subsection{Software Versions}

Table~\ref{tab:software} lists all software dependencies and their versions.

\begin{table}[H]
\centering
\caption{Software dependencies.}
\label{tab:software}
\small
\begin{tabular}{lll}
\toprule
\textbf{Package} & \textbf{Version} & \textbf{Purpose} \\
\midrule
Python           & 3.14             & Runtime \\
PyTorch          & 2.10             & Deep learning framework \\
snnTorch         & 0.9.4            & SNN neuron models and training \\
librosa          & 0.10.x           & Audio loading and mel spectrogram computation \\
NeuroBench       & 2.2.0            & Energy and synaptic operation metrics \\
panns-inference  & latest           & CNN14 embedding extraction \\
torchattacks     & latest           & FGSM and PGD adversarial attacks \\
sPyNNaker        & 1.0.0            & SpiNNaker deployment (Python 3.11 venv) \\
scikit-learn     & latest           & t-SNE, classification metrics \\
scipy            & latest           & Wilcoxon signed-rank test \\
\bottomrule
\end{tabular}
\end{table}

\subsection{Hardware}

\begin{itemize}[nosep]
    \item \textbf{GPU training:} NVIDIA A100-SXM4-80GB on the University of Manchester CSF3 cluster. CUDA 12.6.2. Maximum two GPUs per user.
    \item \textbf{Local development:} Apple M-series with MPS backend (PyTorch). Used for code development and debugging only; all reported results are from CSF3 A100 runs.
    \item \textbf{SpiNNaker:} SpiNN-5 board (48 SpiNNaker chips, 864 ARM968 cores) accessed remotely via \texttt{spinnaker.cs.man.ac.uk}. sPyNNaker 1.0.0.
\end{itemize}

\subsection{Data Preprocessing Pipeline}

The exact preprocessing pipeline, implemented in \texttt{src/dataset.py}, is as follows:

\begin{lstlisting}[language=Python, caption={Audio preprocessing pipeline.}]
y, sr = librosa.load(filepath, sr=22050, duration=5)
expected_len = 22050 * 5  # 110,250 samples
if len(y) < expected_len:
    y = np.pad(y, (0, expected_len - len(y)))
mel = librosa.feature.melspectrogram(
    y=y, sr=sr, n_mels=64, n_fft=1024, hop_length=512
)
mel_db = librosa.power_to_db(mel, ref=np.max)
mel_norm = (mel_db - mel_db.min()) / (mel_db.max() - mel_db.min() + 1e-8)
# Final shape: (1, 64, 216)
tensor = torch.tensor(mel_norm, dtype=torch.float32).unsqueeze(0)
\end{lstlisting}

\subsection{Dataset}

ESC-50 is publicly available at \url{https://github.com/karolpiczak/ESC-50}. It comprises 2{,}000 five-second recordings across 50 classes with predefined five-fold cross-validation splits. No modifications were made to the dataset. Audio files are automatically downloaded and cached at \texttt{data/ESC-50-master/} on first execution.

\subsection{Code Availability}

All source code, training scripts, experiment scripts, analysis notebooks, and result JSON files are available in the project repository. The repository structure is:

\begin{lstlisting}[basicstyle=\small\ttfamily, frame=none]
snn-esc50/
  src/models/        -- SpikingCNN, ConvANN, encoding modules
  src/dataset.py     -- ESC-50 data loading with librosa
  src/train.py       -- 5-fold CV training entry point
  experiments/       -- All advanced experiment scripts
  spinnaker/         -- SpiNNaker deployment scripts
  results/           -- JSON result files (model checkpoints excluded)
\end{lstlisting}
 from report.tex (after \appendix command)

\chapter{Supplementary Material}
\label{ch:appendix}

%=================================================================
\section{Full Per-Fold Results: All Seven Encodings}
\label{sec:app-encoding-results}
%=================================================================

Table~\ref{tab:full-encoding} reports the complete per-fold accuracy for all seven spike encoding methods and the ANN baseline. All experiments use the SpikingCNN architecture ($\sim$622K parameters), ESC-50 five-fold cross-validation (400 test samples per fold), Adam optimiser (lr $= 10^{-3}$, weight decay $= 10^{-4}$), 50 epochs maximum with early stopping at patience $= 10$, and $T = 25$ timesteps. The ANN uses the identical architecture with ReLU activations and no temporal dimension. Population encoding uses 500 output neurons (10 per class) with MSE count loss; all other SNN variants use cross-entropy on summed membrane potentials.

\begin{table}[H]
\centering
\caption{Complete per-fold accuracy (\%) for all encodings on ESC-50 five-fold cross-validation.}
\label{tab:full-encoding}
\small
\begin{tabular}{lccccccr}
\toprule
\textbf{Encoding} & \textbf{Fold 1} & \textbf{Fold 2} & \textbf{Fold 3} & \textbf{Fold 4} & \textbf{Fold 5} & \textbf{Mean} & \textbf{Std} \\
\midrule
ANN (baseline) & 63.25 & 59.50 & 65.25 & 68.75 & 62.50 & 63.85 & 3.07 \\
\midrule
SNN Direct     & 40.50 & 48.50 & 48.25 & 54.00 & 44.50 & 47.15 & 4.50 \\
SNN Rate       & 24.50 & 27.25 & 23.00 & 21.50 & 23.75 & 24.00 & 1.90 \\
SNN Phase      & 22.50 & 22.25 & 25.00 & 24.25 & 26.75 & 24.15 & 1.66 \\
SNN Population & 22.75 & 18.50 & 15.75 & 22.00 & 16.75 & 19.15 & 2.79 \\
SNN Latency    & 14.00 & 15.75 & 17.75 & 15.50 & 18.50 & 16.30 & 1.62 \\
SNN Delta      &  8.25 &  7.75 &  7.25 &  7.50 &  5.50 &  7.25 & 0.94 \\
SNN Burst      &  5.00 &  5.25 &  9.25 &  6.00 &  7.00 &  6.50 & 1.54 \\
\bottomrule
\end{tabular}
\end{table}

%=================================================================
\section{PANNs Transfer Learning: Per-Fold Results}
\label{sec:app-panns}
%=================================================================

Table~\ref{tab:panns-folds} reports the per-fold accuracy for the PANNs transfer learning experiment. Frozen CNN14 embeddings (2,048 dimensions, pretrained on AudioSet) are classified by three heads: a three-layer SNN, a matched three-layer ANN (ReLU), and a linear logistic regression baseline. All heads are trained for 50 epochs with the same Adam optimiser and learning rate schedule.

\begin{table}[H]
\centering
\caption{PANNs transfer learning: per-fold accuracy (\%) on ESC-50.}
\label{tab:panns-folds}
\small
\begin{tabular}{lccccccc}
\toprule
\textbf{Model} & \textbf{Fold 1} & \textbf{Fold 2} & \textbf{Fold 3} & \textbf{Fold 4} & \textbf{Fold 5} & \textbf{Mean} & \textbf{Std} \\
\midrule
PANNs + SNN    & 92.00 & 94.50 & 91.00 & 93.50 & 91.50 & 92.50 & 1.30 \\
PANNs + ANN    & 93.00 & 95.00 & 92.00 & 95.50 & 91.75 & 93.45 & 1.54 \\
PANNs + Linear & 94.25 & 95.75 & 92.50 & 95.25 & 91.25 & 93.80 & 1.69 \\
\bottomrule
\end{tabular}
\end{table}

%=================================================================
\section{SpiNNaker Calibration Parameters}
\label{sec:app-spinnaker-params}
%=================================================================

Table~\ref{tab:spinnaker-neuron} lists the neuron model parameters used for all SpiNNaker deployments. These were determined through a systematic parameter sweep on a 20-sample pilot set (Run~5).

\begin{table}[H]
\centering
\caption{SpiNNaker neuron model parameters (IF\_curr\_exp).}
\label{tab:spinnaker-neuron}
\small
\begin{tabular}{llp{8cm}}
\toprule
\textbf{Parameter} & \textbf{Value} & \textbf{Rationale} \\
\midrule
Neuron model       & IF\_curr\_exp       & sPyNNaker default; approximates snnTorch LIF \\
$v_\text{thresh}$  & 1.0\,mV            & Matches snnTorch threshold $\vartheta = 1.0$ \\
$v_\text{rest}$    & 0.0\,mV            & Matches snnTorch reset-to-zero semantics \\
$v_\text{reset}$   & 0.0\,mV            & Hard reset after spike emission \\
$\tau_m$           & 20.0\,ms           & Calibrated via scale sweep; corresponds to $\beta \approx 0.95$ \\
$\tau_\text{syn,E}$& 5.0\,ms            & Faster decay than $\tau_m$ for temporal precision \\
$\tau_\text{syn,I}$& 5.0\,ms            & Symmetric (no inhibitory synapses used) \\
$C_m$              & 1.0\,nF            & Default capacitance \\
$I_\text{offset}$  & 0.0\,nA            & No bias current injected \\
\bottomrule
\end{tabular}
\end{table}

Table~\ref{tab:spinnaker-board} summarises the SpiNNaker board configuration and simulation parameters.

\begin{table}[H]
\centering
\caption{SpiNNaker board and simulation configuration.}
\label{tab:spinnaker-board}
\small
\begin{tabular}{ll}
\toprule
\textbf{Property} & \textbf{Value} \\
\midrule
Board              & SpiNN-5 (48 chips, 864 ARM cores) \\
Access             & \texttt{spinnaker.cs.man.ac.uk} via sPyNNaker \\
sPyNNaker version  & 1.0.0 \\
Simulation timestep& 1.0\,ms \\
Duration per sample& 25\,ms ($T = 25$ timesteps) \\
FC2 layer size     & 256 input neurons $\to$ 50 output neurons \\
Connectivity       & FromListConnector (pre-computed weight matrix) \\
Recording          & Membrane potential at each timestep for all 50 output neurons \\
\bottomrule
\end{tabular}
\end{table}

%=================================================================
\section{SpiNNaker Pruning Results}
\label{sec:app-spinnaker-pruning}
%=================================================================

Table~\ref{tab:spinnaker-pruning} reports the mean SpiNNaker and snnTorch accuracy across five folds at each pruning level. Pruning is applied as unstructured magnitude pruning to the FC2 weight matrix; weights below the $k$-th percentile are set to zero.

\begin{table}[H]
\centering
\caption{SpiNNaker and snnTorch accuracy (\%) by pruning level (five-fold means).}
\label{tab:spinnaker-pruning}
\small
\begin{tabular}{lccc}
\toprule
\textbf{Pruning (\%)} & \textbf{SpiNNaker} & \textbf{snnTorch} & \textbf{Gap (pp)} \\
\midrule
0  (baseline)  & 33.1 & 46.0 & 12.9 \\
50             & 33.8 & 44.2 & 10.4 \\
55             & 34.1 & 43.5 &  9.4 \\
60             & 35.2 & 42.8 &  7.6 \\
65             & 34.6 & 41.3 &  6.7 \\
70             & 33.9 & 39.7 &  5.8 \\
75             & 32.4 & 37.1 &  4.7 \\
80             & 30.8 & 33.5 &  2.7 \\
85             & 26.1 & 28.9 &  2.8 \\
90             & 19.5 & 21.0 &  1.5 \\
95             &  8.3 &  9.2 &  0.9 \\
\bottomrule
\end{tabular}
\end{table}

The software--hardware gap narrows monotonically with pruning. At 60\% pruning, the SpiNNaker accuracy (35.2\%) is higher than the unpruned software baseline might suggest given the gap trend. At 80\% pruning, the gap has narrowed to 2.7 pp. This supports the interpretation that small-magnitude weights---those most susceptible to fixed-point quantisation error---are disproportionately responsible for the hardware accuracy deficit.

%=================================================================
\section{Reproducibility}
\label{sec:app-reproducibility}
%=================================================================

\subsection{Random Seeds}

Table~\ref{tab:seeds} documents the random seeds used across all experiments to ensure reproducibility.

\begin{table}[H]
\centering
\caption{Random seed configuration for all experiments.}
\label{tab:seeds}
\small
\begin{tabular}{lll}
\toprule
\textbf{Experiment} & \textbf{Seed} & \textbf{Notes} \\
\midrule
All 5-fold training        & 42 (torch, numpy)  & Set at start of each fold \\
CSF3 cluster training      & fold number (1--5) & Per-fold in SLURM submission \\
Surrogate ablation         & 42                 & Fixed single seed, fold 1 \\
t-SNE visualisation        & 42                 & \texttt{sklearn.TSNE random\_state} \\
PANNs + SNN head           & 42                 & All 5 folds \\
Adversarial attacks        & N/A                & Deterministic given model weights \\
Continual learning         & 42                 & Pretrained from fold 4 best \\
\bottomrule
\end{tabular}
\end{table}

\subsection{Software Versions}

Table~\ref{tab:software} lists all software dependencies and their versions.

\begin{table}[H]
\centering
\caption{Software dependencies.}
\label{tab:software}
\small
\begin{tabular}{lll}
\toprule
\textbf{Package} & \textbf{Version} & \textbf{Purpose} \\
\midrule
Python           & 3.14             & Runtime \\
PyTorch          & 2.10             & Deep learning framework \\
snnTorch         & 0.9.4            & SNN neuron models and training \\
librosa          & 0.10.x           & Audio loading and mel spectrogram computation \\
NeuroBench       & 2.2.0            & Energy and synaptic operation metrics \\
panns-inference  & latest           & CNN14 embedding extraction \\
torchattacks     & latest           & FGSM and PGD adversarial attacks \\
sPyNNaker        & 1.0.0            & SpiNNaker deployment (Python 3.11 venv) \\
scikit-learn     & latest           & t-SNE, classification metrics \\
scipy            & latest           & Wilcoxon signed-rank test \\
\bottomrule
\end{tabular}
\end{table}

\subsection{Hardware}

\begin{itemize}[nosep]
    \item \textbf{GPU training:} NVIDIA A100-SXM4-80GB on the University of Manchester CSF3 cluster. CUDA 12.6.2. Maximum two GPUs per user.
    \item \textbf{Local development:} Apple M-series with MPS backend (PyTorch). Used for code development and debugging only; all reported results are from CSF3 A100 runs.
    \item \textbf{SpiNNaker:} SpiNN-5 board (48 SpiNNaker chips, 864 ARM968 cores) accessed remotely via \texttt{spinnaker.cs.man.ac.uk}. sPyNNaker 1.0.0.
\end{itemize}

\subsection{Data Preprocessing Pipeline}

The exact preprocessing pipeline, implemented in \texttt{src/dataset.py}, is as follows:

\begin{lstlisting}[language=Python, caption={Audio preprocessing pipeline.}]
y, sr = librosa.load(filepath, sr=22050, duration=5)
expected_len = 22050 * 5  # 110,250 samples
if len(y) < expected_len:
    y = np.pad(y, (0, expected_len - len(y)))
mel = librosa.feature.melspectrogram(
    y=y, sr=sr, n_mels=64, n_fft=1024, hop_length=512
)
mel_db = librosa.power_to_db(mel, ref=np.max)
mel_norm = (mel_db - mel_db.min()) / (mel_db.max() - mel_db.min() + 1e-8)
# Final shape: (1, 64, 216)
tensor = torch.tensor(mel_norm, dtype=torch.float32).unsqueeze(0)
\end{lstlisting}

\subsection{Dataset}

ESC-50 is publicly available at \url{https://github.com/karolpiczak/ESC-50}. It comprises 2{,}000 five-second recordings across 50 classes with predefined five-fold cross-validation splits. No modifications were made to the dataset. Audio files are automatically downloaded and cached at \texttt{data/ESC-50-master/} on first execution.

\subsection{Code Availability}

All source code, training scripts, experiment scripts, analysis notebooks, and result JSON files are available in the project repository. The repository structure is:

\begin{lstlisting}[basicstyle=\small\ttfamily, frame=none]
snn-esc50/
  src/models/        -- SpikingCNN, ConvANN, encoding modules
  src/dataset.py     -- ESC-50 data loading with librosa
  src/train.py       -- 5-fold CV training entry point
  experiments/       -- All advanced experiment scripts
  spinnaker/         -- SpiNNaker deployment scripts
  results/           -- JSON result files (model checkpoints excluded)
\end{lstlisting}
 from report.tex (after \appendix command)

\chapter{Supplementary Material}
\label{ch:appendix}

%=================================================================
\section{Full Per-Fold Results: All Seven Encodings}
\label{sec:app-encoding-results}
%=================================================================

Table~\ref{tab:full-encoding} reports the complete per-fold accuracy for all seven spike encoding methods and the ANN baseline. All experiments use the SpikingCNN architecture ($\sim$622K parameters), ESC-50 five-fold cross-validation (400 test samples per fold), Adam optimiser (lr $= 10^{-3}$, weight decay $= 10^{-4}$), 50 epochs maximum with early stopping at patience $= 10$, and $T = 25$ timesteps. The ANN uses the identical architecture with ReLU activations and no temporal dimension. Population encoding uses 500 output neurons (10 per class) with MSE count loss; all other SNN variants use cross-entropy on summed membrane potentials.

\begin{table}[H]
\centering
\caption{Complete per-fold accuracy (\%) for all encodings on ESC-50 five-fold cross-validation.}
\label{tab:full-encoding}
\small
\begin{tabular}{lccccccr}
\toprule
\textbf{Encoding} & \textbf{Fold 1} & \textbf{Fold 2} & \textbf{Fold 3} & \textbf{Fold 4} & \textbf{Fold 5} & \textbf{Mean} & \textbf{Std} \\
\midrule
ANN (baseline) & 63.25 & 59.50 & 65.25 & 68.75 & 62.50 & 63.85 & 3.07 \\
\midrule
SNN Direct     & 40.50 & 48.50 & 48.25 & 54.00 & 44.50 & 47.15 & 4.50 \\
SNN Rate       & 24.50 & 27.25 & 23.00 & 21.50 & 23.75 & 24.00 & 1.90 \\
SNN Phase      & 22.50 & 22.25 & 25.00 & 24.25 & 26.75 & 24.15 & 1.66 \\
SNN Population & 22.75 & 18.50 & 15.75 & 22.00 & 16.75 & 19.15 & 2.79 \\
SNN Latency    & 14.00 & 15.75 & 17.75 & 15.50 & 18.50 & 16.30 & 1.62 \\
SNN Delta      &  8.25 &  7.75 &  7.25 &  7.50 &  5.50 &  7.25 & 0.94 \\
SNN Burst      &  5.00 &  5.25 &  9.25 &  6.00 &  7.00 &  6.50 & 1.54 \\
\bottomrule
\end{tabular}
\end{table}

%=================================================================
\section{PANNs Transfer Learning: Per-Fold Results}
\label{sec:app-panns}
%=================================================================

Table~\ref{tab:panns-folds} reports the per-fold accuracy for the PANNs transfer learning experiment. Frozen CNN14 embeddings (2,048 dimensions, pretrained on AudioSet) are classified by three heads: a three-layer SNN, a matched three-layer ANN (ReLU), and a linear logistic regression baseline. All heads are trained for 50 epochs with the same Adam optimiser and learning rate schedule.

\begin{table}[H]
\centering
\caption{PANNs transfer learning: per-fold accuracy (\%) on ESC-50.}
\label{tab:panns-folds}
\small
\begin{tabular}{lccccccc}
\toprule
\textbf{Model} & \textbf{Fold 1} & \textbf{Fold 2} & \textbf{Fold 3} & \textbf{Fold 4} & \textbf{Fold 5} & \textbf{Mean} & \textbf{Std} \\
\midrule
PANNs + SNN    & 92.00 & 94.50 & 91.00 & 93.50 & 91.50 & 92.50 & 1.30 \\
PANNs + ANN    & 93.00 & 95.00 & 92.00 & 95.50 & 91.75 & 93.45 & 1.54 \\
PANNs + Linear & 94.25 & 95.75 & 92.50 & 95.25 & 91.25 & 93.80 & 1.69 \\
\bottomrule
\end{tabular}
\end{table}

%=================================================================
\section{SpiNNaker Calibration Parameters}
\label{sec:app-spinnaker-params}
%=================================================================

Table~\ref{tab:spinnaker-neuron} lists the neuron model parameters used for all SpiNNaker deployments. These were determined through a systematic parameter sweep on a 20-sample pilot set (Run~5).

\begin{table}[H]
\centering
\caption{SpiNNaker neuron model parameters (IF\_curr\_exp).}
\label{tab:spinnaker-neuron}
\small
\begin{tabular}{llp{8cm}}
\toprule
\textbf{Parameter} & \textbf{Value} & \textbf{Rationale} \\
\midrule
Neuron model       & IF\_curr\_exp       & sPyNNaker default; approximates snnTorch LIF \\
$v_\text{thresh}$  & 1.0\,mV            & Matches snnTorch threshold $\vartheta = 1.0$ \\
$v_\text{rest}$    & 0.0\,mV            & Matches snnTorch reset-to-zero semantics \\
$v_\text{reset}$   & 0.0\,mV            & Hard reset after spike emission \\
$\tau_m$           & 20.0\,ms           & Calibrated via scale sweep; corresponds to $\beta \approx 0.95$ \\
$\tau_\text{syn,E}$& 5.0\,ms            & Faster decay than $\tau_m$ for temporal precision \\
$\tau_\text{syn,I}$& 5.0\,ms            & Symmetric (no inhibitory synapses used) \\
$C_m$              & 1.0\,nF            & Default capacitance \\
$I_\text{offset}$  & 0.0\,nA            & No bias current injected \\
\bottomrule
\end{tabular}
\end{table}

Table~\ref{tab:spinnaker-board} summarises the SpiNNaker board configuration and simulation parameters.

\begin{table}[H]
\centering
\caption{SpiNNaker board and simulation configuration.}
\label{tab:spinnaker-board}
\small
\begin{tabular}{ll}
\toprule
\textbf{Property} & \textbf{Value} \\
\midrule
Board              & SpiNN-5 (48 chips, 864 ARM cores) \\
Access             & \texttt{spinnaker.cs.man.ac.uk} via sPyNNaker \\
sPyNNaker version  & 1.0.0 \\
Simulation timestep& 1.0\,ms \\
Duration per sample& 25\,ms ($T = 25$ timesteps) \\
FC2 layer size     & 256 input neurons $\to$ 50 output neurons \\
Connectivity       & FromListConnector (pre-computed weight matrix) \\
Recording          & Membrane potential at each timestep for all 50 output neurons \\
\bottomrule
\end{tabular}
\end{table}

%=================================================================
\section{SpiNNaker Pruning Results}
\label{sec:app-spinnaker-pruning}
%=================================================================

Table~\ref{tab:spinnaker-pruning} reports the mean SpiNNaker and snnTorch accuracy across five folds at each pruning level. Pruning is applied as unstructured magnitude pruning to the FC2 weight matrix; weights below the $k$-th percentile are set to zero.

\begin{table}[H]
\centering
\caption{SpiNNaker and snnTorch accuracy (\%) by pruning level (five-fold means).}
\label{tab:spinnaker-pruning}
\small
\begin{tabular}{lccc}
\toprule
\textbf{Pruning (\%)} & \textbf{SpiNNaker} & \textbf{snnTorch} & \textbf{Gap (pp)} \\
\midrule
0  (baseline)  & 33.1 & 46.0 & 12.9 \\
50             & 33.8 & 44.2 & 10.4 \\
55             & 34.1 & 43.5 &  9.4 \\
60             & 35.2 & 42.8 &  7.6 \\
65             & 34.6 & 41.3 &  6.7 \\
70             & 33.9 & 39.7 &  5.8 \\
75             & 32.4 & 37.1 &  4.7 \\
80             & 30.8 & 33.5 &  2.7 \\
85             & 26.1 & 28.9 &  2.8 \\
90             & 19.5 & 21.0 &  1.5 \\
95             &  8.3 &  9.2 &  0.9 \\
\bottomrule
\end{tabular}
\end{table}

The software--hardware gap narrows monotonically with pruning. At 60\% pruning, the SpiNNaker accuracy (35.2\%) is higher than the unpruned software baseline might suggest given the gap trend. At 80\% pruning, the gap has narrowed to 2.7 pp. This supports the interpretation that small-magnitude weights---those most susceptible to fixed-point quantisation error---are disproportionately responsible for the hardware accuracy deficit.

%=================================================================
\section{Reproducibility}
\label{sec:app-reproducibility}
%=================================================================

\subsection{Random Seeds}

Table~\ref{tab:seeds} documents the random seeds used across all experiments to ensure reproducibility.

\begin{table}[H]
\centering
\caption{Random seed configuration for all experiments.}
\label{tab:seeds}
\small
\begin{tabular}{lll}
\toprule
\textbf{Experiment} & \textbf{Seed} & \textbf{Notes} \\
\midrule
All 5-fold training        & 42 (torch, numpy)  & Set at start of each fold \\
CSF3 cluster training      & fold number (1--5) & Per-fold in SLURM submission \\
Surrogate ablation         & 42                 & Fixed single seed, fold 1 \\
t-SNE visualisation        & 42                 & \texttt{sklearn.TSNE random\_state} \\
PANNs + SNN head           & 42                 & All 5 folds \\
Adversarial attacks        & N/A                & Deterministic given model weights \\
Continual learning         & 42                 & Pretrained from fold 4 best \\
\bottomrule
\end{tabular}
\end{table}

\subsection{Software Versions}

Table~\ref{tab:software} lists all software dependencies and their versions.

\begin{table}[H]
\centering
\caption{Software dependencies.}
\label{tab:software}
\small
\begin{tabular}{lll}
\toprule
\textbf{Package} & \textbf{Version} & \textbf{Purpose} \\
\midrule
Python           & 3.14             & Runtime \\
PyTorch          & 2.10             & Deep learning framework \\
snnTorch         & 0.9.4            & SNN neuron models and training \\
librosa          & 0.10.x           & Audio loading and mel spectrogram computation \\
NeuroBench       & 2.2.0            & Energy and synaptic operation metrics \\
panns-inference  & latest           & CNN14 embedding extraction \\
torchattacks     & latest           & FGSM and PGD adversarial attacks \\
sPyNNaker        & 1.0.0            & SpiNNaker deployment (Python 3.11 venv) \\
scikit-learn     & latest           & t-SNE, classification metrics \\
scipy            & latest           & Wilcoxon signed-rank test \\
\bottomrule
\end{tabular}
\end{table}

\subsection{Hardware}

\begin{itemize}[nosep]
    \item \textbf{GPU training:} NVIDIA A100-SXM4-80GB on the University of Manchester CSF3 cluster. CUDA 12.6.2. Maximum two GPUs per user.
    \item \textbf{Local development:} Apple M-series with MPS backend (PyTorch). Used for code development and debugging only; all reported results are from CSF3 A100 runs.
    \item \textbf{SpiNNaker:} SpiNN-5 board (48 SpiNNaker chips, 864 ARM968 cores) accessed remotely via \texttt{spinnaker.cs.man.ac.uk}. sPyNNaker 1.0.0.
\end{itemize}

\subsection{Data Preprocessing Pipeline}

The exact preprocessing pipeline, implemented in \texttt{src/dataset.py}, is as follows:

\begin{lstlisting}[language=Python, caption={Audio preprocessing pipeline.}]
y, sr = librosa.load(filepath, sr=22050, duration=5)
expected_len = 22050 * 5  # 110,250 samples
if len(y) < expected_len:
    y = np.pad(y, (0, expected_len - len(y)))
mel = librosa.feature.melspectrogram(
    y=y, sr=sr, n_mels=64, n_fft=1024, hop_length=512
)
mel_db = librosa.power_to_db(mel, ref=np.max)
mel_norm = (mel_db - mel_db.min()) / (mel_db.max() - mel_db.min() + 1e-8)
# Final shape: (1, 64, 216)
tensor = torch.tensor(mel_norm, dtype=torch.float32).unsqueeze(0)
\end{lstlisting}

\subsection{Dataset}

ESC-50 is publicly available at \url{https://github.com/karolpiczak/ESC-50}. It comprises 2{,}000 five-second recordings across 50 classes with predefined five-fold cross-validation splits. No modifications were made to the dataset. Audio files are automatically downloaded and cached at \texttt{data/ESC-50-master/} on first execution.

\subsection{Code Availability}

All source code, training scripts, experiment scripts, analysis notebooks, and result JSON files are available in the project repository. The repository structure is:

\begin{lstlisting}[basicstyle=\small\ttfamily, frame=none]
snn-esc50/
  src/models/        -- SpikingCNN, ConvANN, encoding modules
  src/dataset.py     -- ESC-50 data loading with librosa
  src/train.py       -- 5-fold CV training entry point
  experiments/       -- All advanced experiment scripts
  spinnaker/         -- SpiNNaker deployment scripts
  results/           -- JSON result files (model checkpoints excluded)
\end{lstlisting}
